\section{Direcciones de investigación y preguntas abiertas}

\subsection{Replay escalable y priorización}

En tareas de alta dimensión, el replay buffer tradicional no basta para cubrir los nuevos estados que es necesario explorar. La investigación futura debe resolver simultáneamente:

\begin{itemize}
    \item Cómo hacer que el prioritized sampling sea compatible con la estimación TD de múltiples pasos.
    \item Cómo integrar demostraciones y correcciones humanas garantizando al mismo tiempo la diversidad del sampling.
    \item Cómo lograr que el buffer se adapte al desplazamiento de la política (adapt to policy shift), como mediante `exponentially weighted recency`.
\end{itemize}

\begin{importantbox}{Problema abierto: higiene del replay a gran escala}
En el entrenamiento distribuido, cuando cada worker escribe en el buffer hay que garantizar que no haya resampleo duplicado ni muestras omitidas; en el futuro se podría usar un vector clock para rastrear la procedencia de la experiencia.
\end{importantbox}

\subsection{Evaluación robusta y diagnóstico de fallas}

Las fallas del Q-Learning suelen ser encubiertas: reward hacking, caída del retorno, colapso de la política. Además de las métricas de simulación, la investigación también debe resolver:

\begin{itemize}
    \item Extraer una `evidence matrix` y asociarla automáticamente a las acciones.
    \item Construir una `mirrored failure suite` para reproducir transiciones catastróficas anteriores.
    \item Conservar `checkpoints` en producción para su verificación con humano en el circuito (human-in-the-loop).
\end{itemize}

\begin{warningbox}{El punto ciego más común en el diagnóstico}
Los equipos de ingeniería suelen fijarse solo en la curva de recompensa; en cuanto la recompensa se estanca, si eso se interpreta erróneamente como éxito, se pasa por alto el drift de la acción. Se recomienda una vista combinada (recompensa + error TD + tasa de éxito).
\end{warningbox}

\subsection{Orquestación y gobernanza de agentes}

Una vez que el Q-Learning entra en el stack de agentes, la gobernanza debe abarcar varios niveles:

\begin{itemize}
    \item Hacer transparentes las exportaciones de métricas (error TD, Q máxima).
    \item Hacer que el control plane del trabajo de entrenamiento registre (id de checkpoint, semilla de datos, estado de los guardrails).
    \item Establecer un `deployment kill switch` que active un rollback cuando `reward drop > δ`.
\end{itemize}

\begin{knowledgebox}{Lista de verificación de gobernanza}
1. Cada release debe llevar una evidence matrix junto con la línea de tiempo de acciones; 2. el pipeline de entrenamiento debe tratar la `reward variance` como señal de seguridad prioritaria; 3. el proceso de despliegue debe soportar un `fast rollback`.
\end{knowledgebox}

\subsection{Resumen de la sección}

Las direcciones de investigación se concentran en el replay escalable, la evaluación robusta y la gobernanza a nivel de agente; este trabajo determinará la confiabilidad del Q-Learning en tareas más complejas.
La última diapositiva señala las futuras direcciones de investigación: Q-Learning de múltiples pasos, el puente de simulación a realidad (sim-to-real) y la evaluación de seguridad.

\begin{figure}[H]
\centering
\includegraphics[width=0.92\textwidth,page=7]{06_cs224r_qlearning_2025.pdf}
\caption{Diapositiva: direcciones de investigación para un Q-Learning seguro}
\end{figure}
\newpage
\subsection{Transferencia de simulación a realidad y migración de datos}

Antes de desplegar el Q-Learning en un sistema real, suele ser necesaria la transferencia de simulación a realidad (sim-to-real):

\begin{itemize}
    \item Agregar ruido de observación y retardo del actuador en la simulación y volver a entrenar la Q.
    \item Incorporar el checkpoint de simulación al replay buffer del objetivo mediante behavior cloning.
    \item Usar destilación de políticas para que una red más pequeña aproxime la red de simulación.
\end{itemize}

\begin{importantbox}{Puntos clave de supervisión en la transferencia}
1. Confirmar la alineación temporal entre la simulación y la realidad; 2. los datos simulados deben incluir el ruido real de campo; 3. el rollout completo debe quedar registrado para su revisión con humano en el circuito.
\end{importantbox}
